\documentclass[answers,12pt,addpoints]{exam} 
\usepackage{import}

\import{C:/Users/prana/OneDrive/Desktop/MathNotes}{style.tex}

% Header
\newcommand{\name}{Pranav Tikkawar}
\newcommand{\course}{01:XXX:XXX}
\newcommand{\assignment}{Homework n}
\author{\name}
\title{\course \ - \assignment}

\begin{document}
\maketitle

\newpage

\section*{Problem 1}
Prove:
\begin{enumerate}
    \item $\mathbb{R} \times \mathbb{R} \sim \mathbb{R}$.
    \item $\mathbb{R} \setminus \mathbb{N} \sim \mathbb{R}$.
    \item $\mathbb{R} \setminus \mathbb{Q} \sim \mathbb{R}$.
\end{enumerate}

\begin{proof}
We use that $\mathbb{R} \sim \{0,1\}^{\mathbb{N}}$.

\smallskip

(1) Since $\mathbb{R} \sim \{0,1\}^{\mathbb{N}}$, we have
\[
\mathbb{R} \times \mathbb{R}
\sim \{0,1\}^{\mathbb{N}} \times \{0,1\}^{\mathbb{N}}.
\]
Define a map
\[
\Phi : \{0,1\}^{\mathbb{N}} \times \{0,1\}^{\mathbb{N}} \to \{0,1\}^{\mathbb{N}}
\]
by interleaving: for sequences $(a_n)_{n \in \mathbb{N}}$ and $(b_n)_{n \in \mathbb{N}}$, set
\[
c_{2n} = a_n, \quad c_{2n+1} = b_n \quad\text{for all } n \in \mathbb{N},
\]
and define $\Phi((a_n),(b_n)) = (c_k)_{k \in \mathbb{N}}$.

This map is a bijection: given $(c_k)$, we recover $(a_n)$ and $(b_n)$ by
$a_n = c_{2n}$ and $b_n = c_{2n+1}$. Hence
\[
\{0,1\}^{\mathbb{N}} \times \{0,1\}^{\mathbb{N}} \sim \{0,1\}^{\mathbb{N}} \sim \mathbb{R},
\]
so $\mathbb{R} \times \mathbb{R} \sim \mathbb{R}$.

\smallskip

(2) Since $\mathbb{N}$ is countable and $\mathbb{R}$ is uncountable, we clearly have
\[
\mathbb{R} \setminus \mathbb{N} \preceq \mathbb{R}
\]
via the inclusion map. For the other direction, enumerate $\mathbb{N}$ as
$\{n_0,n_1,n_2,\dots\}$ and define a map $f : \mathbb{R} \to \mathbb{R} \setminus \mathbb{N}$ by
\[
f(x) =
\begin{cases}
x, & \text{if } x \notin \mathbb{N},\\[4pt]
k, & \text{if } x = n_k \text{ (so $k \in \mathbb{N}$)}.
\end{cases}
\]
This is well-defined and injective: distinct reals either coincide as non-naturals
where $f$ is the identity, or they are distinct naturals $n_j \neq n_k$ and then
$f(n_j) = j \neq k = f(n_k)$.

Thus we have injections
\[
\mathbb{R} \setminus \mathbb{N} \preceq \mathbb{R}
\quad\text{and}\quad
\mathbb{R} \preceq \mathbb{R} \setminus \mathbb{N},
\]
so by the Cantor--Bernstein--Schröder theorem,
\[
\mathbb{R} \setminus \mathbb{N} \sim \mathbb{R}.
\]

\smallskip

(3) The set $\mathbb{Q}$ is countable, so the same argument applies with
$\mathbb{Q}$ in place of $\mathbb{N}$. Therefore removing a countable subset
from $\mathbb{R}$ does not change its cardinality, and
\[
\mathbb{R} \setminus \mathbb{Q} \sim \mathbb{R}.
\]
\end{proof}

\section*{Problem 2}
Let $A$ be the set of all bijections $f : \mathbb{N} \to \mathbb{N}$. Show that
$A \sim \{0,1\}^{\mathbb{N}}$.

\begin{proof}
We construct injections both ways and apply the Cantor--Bernstein--Schröder theorem.

\smallskip

\emph{Injection $\{0,1\}^{\mathbb{N}} \to A$.}
Given $x = (x_n)_{n\in\mathbb{N}} \in \{0,1\}^{\mathbb{N}}$, define $f_x : \mathbb{N} \to \mathbb{N}$ as follows.
For each $n \in \mathbb{N}$, consider the pair $\{2n, 2n+1\}$:
\begin{itemize}
    \item If $x_n = 0$, let $f_x(2n) = 2n$ and $f_x(2n+1) = 2n+1$.
    \item If $x_n = 1$, let $f_x(2n) = 2n+1$ and $f_x(2n+1) = 2n$.
\end{itemize}
This defines a permutation of $\mathbb{N}$, hence a bijection, so $f_x \in A$.
If $x \neq y$, there is some $n$ with $x_n \neq y_n$, and on the pair
$\{2n,2n+1\}$ the maps $f_x$ and $f_y$ act differently. Thus $f_x \neq f_y$.
Therefore
\[
\Phi : \{0,1\}^{\mathbb{N}} \to A,\quad \Phi(x) = f_x
\]
is injective.

\smallskip

\emph{Injection $A \to \{0,1\}^{\mathbb{N}}$.}
We know from class that $\mathbb{N}^{\mathbb{N}} \sim \{0,1\}^{\mathbb{N}}$.
Since $A \subseteq \mathbb{N}^{\mathbb{N}}$, the inclusion map
\[
\iota : A \hookrightarrow \mathbb{N}^{\mathbb{N}}
\]
is injective. Composing $\iota$ with a fixed bijection
$\psi : \mathbb{N}^{\mathbb{N}} \to \{0,1\}^{\mathbb{N}}$ gives an injection
\[
\psi \circ \iota : A \to \{0,1\}^{\mathbb{N}}.
\]

We now have
\[
\{0,1\}^{\mathbb{N}} \preceq A
\quad\text{and}\quad
A \preceq \{0,1\}^{\mathbb{N}},
\]
so by Cantor--Bernstein--Schröder,
\[
A \sim \{0,1\}^{\mathbb{N}}.
\]
\end{proof}

\section*{Problem 3}
For sets $A,B$ let $A \Delta B = (A \setminus B) \cup (B \setminus A)$.
\begin{enumerate}
    \item[(i)] Show $A = B$ iff $A \Delta B = \varnothing$.
    \item[(ii)] Show $(A \Delta B) \cup (B \Delta C) \subseteq A \Delta C$.
    \item[(iii)] Let
    \[
    R = \{\langle A,B\rangle \in \mathcal{P}(\mathbb{N})^2 : 2026 \notin A \Delta B\}.
    \]
    Show $R$ is an equivalence relation on $\mathcal{P}(\mathbb{N})$.
    \item[(iv)] Compute $[\varnothing]_R$, $[\mathbb{N}]_R$ and the number of
    equivalence classes.
\end{enumerate}

\begin{proof}
(i) Suppose $A = B$. Then $A \setminus B = \varnothing$ and $B \setminus A = \varnothing$,
so $A \Delta B = \varnothing \cup \varnothing = \varnothing$.

Conversely, if $A \Delta B = \varnothing$, then no element lies in exactly one of $A,B$.
Hence every element of $A$ lies in $B$ and every element of $B$ lies in $A$, so $A=B$.

\smallskip

(ii) For any set $X$ and element $x$, let $1_X(x)$ be $1$ if $x \in X$ and $0$ otherwise.
Then $x \in A \Delta B$ iff $1_A(x) \oplus 1_B(x) = 1$, where $\oplus$ is addition modulo $2$.
Similarly, $x \in B \Delta C$ iff $1_B(x) \oplus 1_C(x) = 1$, and
$x \in A \Delta C$ iff $1_A(x) \oplus 1_C(x) = 1$.

Suppose $x \in (A \Delta B) \cup (B \Delta C)$. Then at least one of
$1_A(x) \oplus 1_B(x)$ and $1_B(x) \oplus 1_C(x)$ equals $1$.
But in $\{0,1\}$ we have
\[
(1_A(x) \oplus 1_B(x)) \oplus (1_B(x) \oplus 1_C(x))
= 1_A(x) \oplus 1_C(x).
\]
If at least one of the two terms on the left is $1$, then the right-hand side is $1$,
so $1_A(x) \oplus 1_C(x) = 1$. Hence $x \in A \Delta C$.
Therefore $(A \Delta B) \cup (B \Delta C) \subseteq A \Delta C$.

\smallskip

(iii) We verify that $R$ is an equivalence relation on $\mathcal{P}(\mathbb{N})$.

\emph{Reflexive:} For every $A \subseteq \mathbb{N}$ we have $A \Delta A = \varnothing$,
so in particular $2026 \notin A \Delta A$. Hence $\langle A,A\rangle \in R$.

\emph{Symmetric:} If $\langle A,B\rangle \in R$, then $2026 \notin A \Delta B$.
But $A \Delta B = B \Delta A$, so $2026 \notin B \Delta A$, and thus $\langle B,A\rangle \in R$.

\emph{Transitive:} Suppose $\langle A,B\rangle \in R$ and $\langle B,C\rangle \in R$.
Then $2026 \notin A \Delta B$ and $2026 \notin B \Delta C$, which means
$2026$ is in both $A$ and $B$ or in neither, and also in both $B$ and $C$ or in neither.
Hence $2026$ is in $A$ and $C$ together or in neither, so $2026 \notin A \Delta C$,
and $\langle A,C\rangle \in R$.

Thus $R$ is an equivalence relation.

\smallskip

(iv) By definition, $A$ and $B$ are $R$-equivalent iff $2026$ is either in both
$A,B$ or in neither. Therefore
\[
[\varnothing]_R = \{A \subseteq \mathbb{N} : 2026 \notin A\},
\]
and
\[
[\mathbb{N}]_R = \{A \subseteq \mathbb{N} : 2026 \in A\}.
\]
There are exactly two equivalence classes: all subsets of $\mathbb{N}$ that do not contain
$2026$, and all subsets that do contain $2026$.
\end{proof}

\section*{Problem 4 (Optional)}
On $\{0,1\}^{\mathbb{N}}$, define $f \sim g$ iff $f(0) = g(0)$, $f(1) = g(1)$, and $f(2) = g(2)$.
How many different equivalence classes are there?

\begin{proof}
Two functions $f,g \in \{0,1\}^{\mathbb{N}}$ are equivalent iff they agree on the first
three coordinates. Thus each equivalence class is determined by the triple
\[
(f(0), f(1), f(2)) \in \{0,1\}^3.
\]
There are $2^3 = 8$ possible such triples, hence there are $8$ equivalence classes.
\end{proof}

\section*{Problem 5}
Let $f : A \to B$ and $g : B \to A$ be injections, and let $h = g \circ f : A \to A$.
We construct a bijection $k : A \to \operatorname{rng}(g)$ and then $g^{-1} \circ k : A \to B$,
obtaining the Cantor--Bernstein--Schröder theorem. As an application, we get a bijection
between $[0,1)$ and $[0,1]$.

\begin{proof}
(I) Since $f$ and $g$ are injections, their composition $h = g \circ f$ is also injective.

(II) The map $g : B \to \operatorname{rng}(g)$ is a bijection by definition of the range,
hence it has an inverse $g^{-1} : \operatorname{rng}(g) \to B$. Thus, if we construct
a bijection $k : A \to \operatorname{rng}(g)$, then
\[
g^{-1} \circ k : A \to B
\]
is a bijection as the composition of bijections.

(III) Define subsets $A_0, A_1, \dots$ of $A$ recursively by
\[
A_0 = A \setminus \operatorname{rng}(g), \qquad A_{n+1} = h[A_n] \quad\text{for all } n \in \mathbb{N}.
\]
Since $h$ is injective, the restriction $h \upharpoonright A_n : A_n \to A_{n+1}$ is a bijection
from $A_n$ onto $A_{n+1}$ for each $n$. The sets $A_n$ are pairwise disjoint by construction.

(IV) Define $k : A \to \operatorname{rng}(g)$ by
\[
k(a) =
\begin{cases}
h(a), & \text{if } a \in A_n \text{ for some } n,\\[4pt]
a,    & \text{if } a \notin \bigcup_{n \in \mathbb{N}} A_n.
\end{cases}
\]
If $a \in A_n$, then $h(a) \in A_{n+1} \subseteq A \subseteq \operatorname{rng}(g) \cup A_0$,
and by construction $A_0 \subseteq A \setminus \operatorname{rng}(g)$, so in fact $h(a) \in \operatorname{rng}(g)$.
If $a \notin \bigcup_n A_n$, then by definition $a \in \operatorname{rng}(g)$. Thus
$k(a) \in \operatorname{rng}(g)$ in either case, so $k$ is well-defined with codomain $\operatorname{rng}(g)$.

(V) To see that $k$ is injective, note first that
\begin{quote}
for any $a \in A$, $k(a) \in A_{n+1}$ for some $n$ iff $a \in A_n$ for some $n$.
\end{quote}
This uses that $h[A_n] = A_{n+1}$ and that the $A_n$ are pairwise disjoint. Now suppose
$k(a) = k(a') = b$. If $b \in A_{m+1}$ for some $m$, then by the claim both $a,a'$ lie in $A_m$,
and on $A_m$ the map $k$ coincides with $h$, which is injective, so $a=a'$.
If $b \notin A_n$ for any $n$, then $b$ is fixed by $k$, so $a = b = a'$.
Thus $k$ is injective.

(VI) To see that $k$ is surjective onto $\operatorname{rng}(g)$, let $y \in \operatorname{rng}(g)$.
If $y \notin A_{n+1}$ for any $n$, then by definition $k(y) = y$, so $y$ is in the range of $k$.
If $y \in A_{n+1}$ for some $n$, then there is a unique $x \in A_n$ with $h(x) = y$,
because $h \upharpoonright A_n$ is a bijection onto $A_{n+1}$. For this $x$ we have
$k(x) = h(x) = y$, so again $y$ is in the range of $k$. Therefore $k : A \to \operatorname{rng}(g)$
is a bijection, and thus $g^{-1} \circ k : A \to B$ is a bijection.

This proves the Cantor--Bernstein--Schröder theorem.

\medskip

\noindent\textbf{Application: a bijection $[0,1) \to [0,1]$.}

Define injections
\[
f_0 : [0,1) \to [0,1], \quad f_0(x) = x,
\]
and
\[
g_0 : [0,1] \to [0,1), \quad g_0(x) = \frac{x}{2}.
\]
Both $f_0$ and $g_0$ are injective. By the Cantor--Bernstein--Schröder theorem,
there exists a bijection between $[0,1)$ and $[0,1]$.
\end{proof}


\end{document}